%%% FIELD proof-prop-no-breakout-nonvacuity
Put \(V_e=\epsilon_{u,e}\) and \(W_e=\epsilon_{t,e}\).  We first compute the environmental cumulant profiles of the witness.  If \(N\sim\operatorname{Poisson}(\lambda)\), then, directly from the Poisson moment generating function,
\[
 \log \mathbb E\exp\{t j(N-\lambda)\}
 =\lambda\bigl(\exp(jt)-1-jt\bigr).
\]
By the definition of a cumulant as a derivative at zero of the log moment generating function, for every integer \(k\geq2\),
\[
 \kappa_k\bigl(j(N-\lambda)\bigr)=\lambda j^k.                 \tag{1}
\]
For fixed \(e\), independence of the variables \(N_{e,j}\) makes the moment generating function of their sum a product, so its logarithm is the sum of their log moment generating functions.  Thus (1), the witness definition, and the definition of \(u_k\) give
\[
 \begin{aligned}
 (u_k)_e
 &=\kappa_k(V_e)-\kappa_k(V_0)\\
 &=\sum_{j=1}^d(\lambda_{e,j}-\lambda_{0,j})j^k\\
 &=\begin{cases}
      \eta e^k,&1\leq e\leq d,\\
      0,&d<e\leq E-1.
    \end{cases}
 \end{aligned}                                                   \tag{2}
\]
The final case is absent when \(E=d+1\).  Consequently,
\[
 u_k=\eta(1^k,2^k,\ldots,d^k,0,\ldots,0)^\top.                 \tag{3}
\]

The treatment equation in this witness is \(T_e=V_e+W_e\).  By source independence, the cumulant of this sum is the sum of the source cumulants, by the same log-moment-generating-function calculation used above.  The variables \(W_e\) have one common Gaussian law in all environments, so their cumulant contrasts are zero.  It follows from the definition of \(c_k\) and (3) that
\[
 c_k=u_k\qquad(k\geq2).                                         \tag{4}
\]
In particular, using the definition \(C=[c_2\mid\cdots\mid c_{d+1}]\),
\[
 C=\begin{bmatrix}\eta V_d\\[2pt]0\end{bmatrix},
 \qquad
 V_d=[j^k]_{j=1,\ldots,d;\ k=2,\ldots,d+1}.                    \tag{5}
\]

We verify the asserted rank rather than assuming it.  Factoring \(j^2\) from row \(j\) of \(V_d\) leaves the matrix \([j^\ell]_{j=1,\ldots,d;\ \ell=0,\ldots,d-1}\).  If a linear combination of its columns were zero, the associated polynomial of degree at most \(d-1\) would vanish at the \(d\) distinct points \(1,\ldots,d\).  Repeated application of the factor theorem would then make it divisible by a polynomial of degree \(d\), so it must be the zero polynomial.  Its columns, and hence those of \(V_d\), are therefore linearly independent.  Since \(\eta>0\), (5) yields
\[
 C^\top C=\eta^2V_d^\top V_d,
 \qquad
 \sigma_d(C)=\eta\sigma_d(V_d)\geq\kappa_0+\rho.              \tag{6}
\]
Thus \(\operatorname{rank}(C)=d\), so the projector \(Q_A\) is well defined.  Equations (3)--(5) also show
\[
 A=\operatorname{col}(C)
  =\operatorname{span}\{\mathbf e_1,\ldots,\mathbf e_d\},
 \qquad
 c_k,u_k\in A\quad(k\geq2).                                   \tag{7}
\]
By the definition of an orthogonal projector, (7) implies \(Q_Ac_k=0\) for every \(k\geq2\).  The definition of the first breakout therefore gives
\[
 r^\star(\theta)=\infty.                                       \tag{8}
\]

We next check membership in the bounded interior.  Every constructed source is centered: this is immediate for \(N_{e,j}-\lambda_{e,j}\), and the Gaussian sources have mean zero.  Mutual independence within each environment holds by the witness construction, as required by source independence.  The displayed Poisson log moment generating function is finite for both signs of every real argument, as are the Gaussian moment generating functions.  For any such source \(X\), any \(h>0\), and every positive integer \(s\),
\[
 |X|^s\leq s!h^{-s}\exp(h|X|),
 \qquad
 \exp(h|X|)\leq\exp(hX)+\exp(-hX).
\]
Taking expectations proves \(\mathbb E|X|^s<\infty\), which verifies the all-absolute-moments member property.  The feasibility inequality in the witness definition gives the required \((2R_{\max}+\delta)\)-moment bound for \(V_e\) and \(\epsilon_y\).  It also gives that bound for every \(W_e\), because all \(W_e\) have the same law as \(W_0\).  Thus the moment envelope holds with \(\rho\) slack.

For the coefficients, the witness and its feasibility conditions give, separately,
\[
 |\beta|=q/2\leq B-\rho,
 \qquad
 a_0+\rho\leq|\alpha|=1\leq a_1-\rho,
 \qquad
 q_0+\rho\leq|\gamma/\alpha|=q\leq q_1-\rho.                  \tag{9}
\]
The positive-profile and environment-indexing conditions hold by the standing assumptions and \(E\geq d+1\); the structural equations hold by construction.  Rank and the calibration gap follow from (6), while the latent-profile-plane property follows from (7).  Together with the centeredness, independence, all-order moments, and moment bound just checked, these are all member properties of \(\Theta\), and (6) and (9) are precisely the additional slack inequalities in the definition of \(\Theta_\rho^\circ\).  Hence
\[
 \theta\in\Theta_\rho^\circ.                                \tag{10}
\]

It remains to check the source swap.  In the original representation, direct substitution of \(\alpha=1\), \(\beta=-q/2\), and \(\gamma=q\) gives, for every \(e\in\mathcal E\),
\[
 T_e=V_e+W_e,
 \qquad
 Y_e=-\frac q2(V_e+W_e)+qV_e+\epsilon_y
     =\frac q2V_e-\frac q2W_e+\epsilon_y.                       \tag{11}
\]
By the definition of \(S\), the swapped structural parameter has
\[
 \alpha'=1,
 \qquad
 \beta'=\beta+\gamma/\alpha=q/2,
 \qquad
 \gamma'=-\gamma/\alpha=-q,
 \qquad
 \epsilon'_{u,e}=W_e,
 \qquad
 \epsilon'_{t,e}=V_e.                                          \tag{12}
\]
Its structural equations consequently produce
\[
 T'_e=W_e+V_e=T_e
\]
and, using (12),
\[
 Y'_e=\frac q2(W_e+V_e)-qW_e+\epsilon_y
     =\frac q2V_e-\frac q2W_e+\epsilon_y
     =Y_e.                                                       \tag{13}
\]
Therefore the two parameters induce exactly the same observed law \(P_e=\mathcal L(T_e,Y_e)\) in every environment, rather than merely matching finitely many moments.  In particular, their observed contrasts \(c_k\), their matrix \(C\), their plane \(A\), and their projector \(Q_A\) coincide.

The swapped latent source \(W_e\) has the same Gaussian law in every environment.  Hence its latent cumulant-shift profile is zero at every order and so belongs to \(A\).  By (4), the unchanged observed treatment contrasts all lie in \(A\), whence
\[
 r^\star(S\theta)=\infty.                                      \tag{14}
\]
Centeredness, source independence, and all-order moment existence are preserved when the two independent source families are interchanged.  Because \(\alpha=1\), the source-level moment maximum is also unchanged by this interchange.  The swapped effect, loading, and confounding magnitudes are, by (12), \(q/2\), \(1\), and \(q\), respectively, so the same three inequalities (9) apply.  Finally, the observed \(C\) is unchanged, so (6) supplies the same rank and singular-value gap.  These checks establish
\[
 S\theta\in\Theta_\rho^\circ.                                \tag{15}
\]

Combining (8), (10), (14), and (15) with the definition of \(\mathcal N_\rho\) yields \(\theta\in\mathcal N_\rho\).  The two observationally identical structural parameters have effects \(-q/2\) and \(q/2\), which are unequal because \(q>0\).  Thus \(\mathcal N_\rho\neq\varnothing\) for every collection of constants satisfying the feasibility conditions in the proposition.
